% GENERATED FILE. Edit canonical lessons, then run npm run generate:books.
% canonical-source-hash: fnv1a64:4d1a287c4046997f
% canonical-lessons: BN-C515-poribeshon, BN-C515-gras, BN-C515-svad, BN-C515-eto, BN-C515-achano, BN-R515-review-508-515

\chapter{At the Table}
\label{ch:bn-at-the-table}
\hlchaptermodality{515}

\begin{chapteropening}
\textbf{By the end of this chapter:} \emph{I can talk about serving food, mouthfuls and leftovers.}

Complete the last lesson of chapter 515: I can talk about serving food, mouthfuls and leftovers.
\end{chapteropening}

\section[\texorpdfstring{\bn{পরিবেশন} (pôribeshôn)}{পরিবেশন (pôribeshôn)}]{\bn{পরিবেশন} (pôribeshôn) — serving food}

\label{lesson:BN-C515-poribeshon}



\begin{quote}
Before the new one: say the Bengali for to steam, then the Bengali for to sauté.
\end{quote}

\subsection*{You'll want to know: \bn{পরিবেশন}}
\textbf{\bn{পরিবেশন}} — \emph{pôribeshôn} — \textquotedblleft{}serving food\textquotedblright{}.

\begin{quote}
\bn{বিয়েতে} \bn{ওয়েটাররা} \bn{খাবার} \bn{পরিবেশন} \bn{করল।}
\end{quote}

\begin{quote}
\emph{Biyete oyeṭārrā khābār pôribeshôn kôrlô.}
\end{quote}

\begin{quote}
At the wedding the waiters served the food.
\end{quote}

\begin{cousinweb}
\textbf{\bn{পরিবেশন}} (\emph{pôribeshôn}) is a tatsama, Sanskrit \emph{pariveṣaṇa}, \textquotedblleft{}waiting at table, serving up a meal\textquotedblright{} (Monier-Williams).
\end{cousinweb}

\begin{grammarlens}[title={What you've built}]
One more word for describing what you lived and felt.
\end{grammarlens}

\subsection*{Your turn}
\begin{itemize}
\raggedright
  \item \emph{Say it:} \emph{pôribeshôn}
  \item \emph{Say it:} \emph{pôribeshôn}, once more
  \item \emph{Say it:} \emph{Biyete oyeṭārrā khābār pôribeshôn kôrlô.} — At the wedding the waiters served the food.
  \item \emph{Recall:} say the Bengali for to steam, then the Bengali for to sauté, then say \emph{pôribeshôn} again
\end{itemize}

\begin{tcolorbox}[breakable,colback=teal!4,colframe=teal!35!black,title={Before you move on}]
What does \bn{পরিবেশন} mean? (\textquotedblleft{}Serving food\textquotedblright{}.) Say it once more.
\end{tcolorbox}
\section[\texorpdfstring{\bn{গ্রাস} (grās)}{গ্রাস (grās)}]{\bn{গ্রাস} (grās) — a mouthful}

\label{lesson:BN-C515-gras}



\begin{quote}
Before the new one: say the Bengali for to sauté, then the Bengali for serving food.
\end{quote}

\subsection*{You'll want to know: \bn{গ্রাস}}
\textbf{\bn{গ্রাস}} — \emph{grās} — \textquotedblleft{}a mouthful\textquotedblright{}.

\begin{quote}
\bn{শিশুটা} \bn{দুই} \bn{গ্রাস} \bn{ভাত} \bn{খেল।}
\end{quote}

\begin{quote}
\emph{Shishuṭā dui grās bhāt khelô.}
\end{quote}

\begin{quote}
The child ate two mouthfuls of rice.
\end{quote}

\begin{cousinweb}
\textbf{\bn{গ্রাস}} (\emph{grās}) is a tatsama, Sanskrit \emph{grāsa}, \textquotedblleft{}a mouthful, a morsel\textquotedblright{} (Monier-Williams).
\end{cousinweb}

\begin{grammarlens}[title={What you've built}]
One more word for describing what you lived and felt.
\end{grammarlens}

\subsection*{Your turn}
\begin{itemize}
\raggedright
  \item \emph{Say it:} \emph{grās}
  \item \emph{Say it:} \emph{grās}, once more
  \item \emph{Say it:} \emph{Shishuṭā dui grās bhāt khelô.} — The child ate two mouthfuls of rice.
  \item \emph{Recall:} say the Bengali for to sauté, then the Bengali for serving food, then say \emph{grās} again
\end{itemize}

\begin{tcolorbox}[breakable,colback=teal!4,colframe=teal!35!black,title={Before you move on}]
What does \bn{গ্রাস} mean? (\textquotedblleft{}A mouthful\textquotedblright{}.) Say it once more.
\end{tcolorbox}
\section[\texorpdfstring{\bn{স্বাদ} (svād)}{স্বাদ (svād)}]{\bn{স্বাদ} (svād) — a taste; a flavour}

\label{lesson:BN-C515-svad}



\begin{quote}
Before the new one: say the Bengali for serving food, then the Bengali for a mouthful.
\end{quote}

\subsection*{You'll want to know: \bn{স্বাদ}}
\textbf{\bn{স্বাদ}} — \emph{svād} — \textquotedblleft{}a taste; a flavour\textquotedblright{}.

\begin{quote}
\bn{এই} \bn{তরকারির} \bn{স্বাদ} \bn{ভালো।}
\end{quote}

\begin{quote}
\emph{Ei tôrkārir svād bhālo.}
\end{quote}

\begin{quote}
This curry tastes good.
\end{quote}

\begin{cousinweb}
\textbf{\bn{স্বাদ}} (\emph{svād}) is a tatsama, Sanskrit \emph{svāda}, \textquotedblleft{}taste, flavour\textquotedblright{} (Monier-Williams).
\end{cousinweb}

\begin{grammarlens}[title={What you've built}]
One more word for describing what you lived and felt.
\end{grammarlens}

\subsection*{Your turn}
\begin{itemize}
\raggedright
  \item \emph{Say it:} \emph{svād}
  \item \emph{Say it:} \emph{svād}, once more
  \item \emph{Say it:} \emph{Ei tôrkārir svād bhālo.} — This curry tastes good.
  \item \emph{Recall:} say the Bengali for serving food, then the Bengali for a mouthful, then say \emph{svād} again
\end{itemize}

\begin{tcolorbox}[breakable,colback=teal!4,colframe=teal!35!black,title={Before you move on}]
What does \bn{স্বাদ} mean? (\textquotedblleft{}A taste; a flavour\textquotedblright{}.) Say it once more.
\end{tcolorbox}
\section[\texorpdfstring{\bn{এঁটো} (ẽṭo)}{এঁটো (ẽṭo)}]{\bn{এঁটো} (ẽṭo) — leftover from a meal; touched by eating}

\label{lesson:BN-C515-eto}



\begin{quote}
Before the new one: say the Bengali for a mouthful, then the Bengali for a flavour.
\end{quote}

\subsection*{You'll want to know: \bn{এঁটো}}
\textbf{\bn{এঁটো}} — \emph{ẽṭo} — \textquotedblleft{}leftover from a meal; touched by eating\textquotedblright{}.

\begin{quote}
\bn{খাওয়ার} \bn{পরে} \bn{আমি} \bn{এঁটো} \bn{থালাগুলো} \bn{বেসিনে} \bn{রাখলাম।}
\end{quote}

\begin{quote}
\emph{Khāoyār pôre āmi ẽṭo thālāgulo besine rākhlām.}
\end{quote}

\begin{quote}
After the meal I put the used plates in the basin.
\end{quote}

\begin{cousinweb}
\textbf{\bn{এঁটো}} (\emph{ẽṭo}) comes from a reconstructed \emph{ācaṣṭa}, \textquotedblleft{}tasted\textquotedblright{}; CDIAL lists Bengali \emph{ẽṭo}, \textquotedblleft{}the leavings of a meal\textquotedblright{} (CDIAL 1067). Food, or a plate, that someone has already eaten from is \bn{এঁটো}.
\end{cousinweb}

\begin{grammarlens}[title={What you've built}]
One more word for describing what you lived and felt.
\end{grammarlens}

\subsection*{Your turn}
\begin{itemize}
\raggedright
  \item \emph{Say it:} \emph{ẽṭo}
  \item \emph{Say it:} \emph{ẽṭo}, once more
  \item \emph{Say it:} \emph{Khāoyār pôre āmi ẽṭo thālāgulo besine rākhlām.} — After the meal I put the used plates in the basin.
  \item \emph{Recall:} say the Bengali for a mouthful, then the Bengali for a flavour, then say \emph{ẽṭo} again
\end{itemize}

\begin{tcolorbox}[breakable,colback=teal!4,colframe=teal!35!black,title={Before you move on}]
What does \bn{এঁটো} mean? (\textquotedblleft{}Leftover from a meal; touched by eating\textquotedblright{}.) Say it once more.
\end{tcolorbox}
\section[\texorpdfstring{\bn{আঁচানো} (ā̃chāno)}{আঁচানো (ā̃chāno)}]{\bn{আঁচানো} (ā̃chāno) — to rinse after a meal (the mouth and hands)}

\label{lesson:BN-C515-achano}



\begin{quote}
Before the new one: say the Bengali for a flavour, then the Bengali for leftover from a meal.
\end{quote}

\subsection*{You'll want to know: \bn{আঁচানো}}
\textbf{\bn{আঁচানো}} — \emph{ā̃chāno} — \textquotedblleft{}to rinse after a meal (the mouth and hands)\textquotedblright{}.

\begin{quote}
\bn{খাওয়ার} \bn{পরে} \bn{সবাই} \bn{উঠে} \bn{আঁচাল।}
\end{quote}

\begin{quote}
\emph{Khāoyār pôre sôbāi uṭhe ā̃chālô.}
\end{quote}

\begin{quote}
After the meal everyone got up and rinsed their mouths and hands.
\end{quote}

Past (he/she …): \bn{তিনি} \bn{খাওয়ার} \bn{পরে} \bn{আঁচালেন}, \emph{tini khāoyār pôre ā̃chālen}.

Rice is often eaten by hand, so mouth and hand are rinsed when the meal is over.

\begin{cousinweb}
\textbf{\bn{আঁচানো}} (\emph{ā̃chāno}) comes down from Sanskrit \emph{ācāmati}, taking a sip of water from the hollow of the hand (CDIAL 1069, which lists Bengali \emph{ā̃cāna}). The everyday sense, rinsing mouth and hand at the end of a meal, probably grew out of that sip.
\end{cousinweb}

\begin{grammarlens}[title={What you've built}]
One more word for describing what you lived and felt.
\end{grammarlens}

\subsection*{Your turn}
\begin{itemize}
\raggedright
  \item \emph{Say it:} \emph{ā̃chāno}
  \item \emph{Say it:} \emph{ā̃chāno}, once more
  \item \emph{Say it:} \emph{Khāoyār pôre sôbāi uṭhe ā̃chālô.} — After the meal everyone got up and rinsed their mouths and hands.
  \item \emph{Recall:} say the Bengali for a flavour, then the Bengali for leftover from a meal, then say \emph{ā̃chāno} again
\end{itemize}

\begin{tcolorbox}[breakable,colback=teal!4,colframe=teal!35!black,title={Before you move on}]
What does \bn{আঁচানো} mean? (\textquotedblleft{}To rinse after a meal (the mouth and hands)\textquotedblright{}.) Say it once more.
\end{tcolorbox}
\section[(dialogue)]{Looking back: from Rooms of a House to At the Table}

\label{lesson:BN-R515-review-508-515}



\begin{quote}
Say the Bengali for a drawing room and an attic.
\end{quote}

\subsection*{Your turn}
Say these aloud:

\begin{itemize}
\raggedright
  \item a drawing room, an attic, a kitchen, a door frame, to hang up
  \item a rug, an embroidered quilt, a doormat, a cushion, to spread out
  \item floor disinfectant, cobwebs, cleared away, to wring out, to scour
  \item dishes, a wide cooking pot, a container, to pour, to dip
  \item a stove, a spatula, tongs, a lid, to grind
  \item to chop, a peel, to grate, to knead, to chop finely
  \item to bring to the boil, to heat up, to toast, to steam, to sauté
  \item serving food, a mouthful, a flavour, leftover from a meal, to rinse after a meal
\end{itemize}

\begin{tcolorbox}[breakable,colback=teal!4,colframe=teal!35!black,title={Before you move on}]
A drawing room? (\textbf{\bn{বৈঠকখানা}}, \emph{boiṭhôkkhānā}.) An attic? (\textbf{\bn{চিলেকোঠা}}, \emph{chilekoṭhā}.) A kitchen? (\textbf{\bn{রান্নাঘর}}, \emph{rānnāghôr}.) A door frame? (\textbf{\bn{চৌকাঠ}}, \emph{choukāṭh}.) To hang up? (\textbf{\bn{টাঙানো}}, \emph{ṭāngāno}.) A rug? (\textbf{\bn{গালিচা}}, \emph{gālichā}.) An embroidered quilt? (\textbf{\bn{কাঁথা}}, \emph{kā̃thā}.) A doormat? (\textbf{\bn{পাপোশ}}, \emph{pāposh}.) A cushion? (\textbf{\bn{গদি}}, \emph{gôdi}.) To spread out? (\textbf{\bn{বিছানো}}, \emph{bichhāno}.) Floor disinfectant? (\textbf{\bn{ফিনাইল}}, \emph{phināil}.) Cobwebs? (\textbf{\bn{ঝুল}}, \emph{jhul}.) Cleared away? (\textbf{\bn{সাফ}}, \emph{sāph}.) To wring out? (\textbf{\bn{নিংড়ানো}}, \emph{ningṛāno}.) To scour? (\textbf{\bn{মাজা}}, \emph{mājā}.) Dishes? (\textbf{\bn{বাসন}}, \emph{bāsôn}.) A wide cooking pot? (\textbf{\bn{ডেকচি}}, \emph{ḍekchi}.) A container? (\textbf{\bn{পাত্র}}, \emph{pātrô}.) To pour? (\textbf{\bn{ঢালা}}, \emph{ḍhālā}.) To dip? (\textbf{\bn{ডোবানো}}, \emph{ḍobāno}.) A stove? (\textbf{\bn{উনুন}}, \emph{unun}.) A spatula? (\textbf{\bn{খুন্তি}}, \emph{khunti}.) Tongs? (\textbf{\bn{সাঁড়াশি}}, \emph{sā̃ṛāshi}.) A lid? (\textbf{\bn{ঢাকনা}}, \emph{ḍhāknā}.) To grind? (\textbf{\bn{বাটা}}, \emph{bāṭā}.) To chop? (\textbf{\bn{কোটা}}, \emph{koṭā}.) A peel? (\textbf{\bn{খোসা}}, \emph{khosā}.) To grate? (\textbf{\bn{কোরা}}, \emph{korā}.) To knead? (\textbf{\bn{মাখা}}, \emph{mākhā}.) To chop finely? (\textbf{\bn{কুচোনো}}, \emph{kuchono}.) To bring to the boil? (\textbf{\bn{ফোটানো}}, \emph{phoṭāno}.) To heat up? (\textbf{\bn{তাতানো}}, \emph{tātāno}.) To toast? (\textbf{\bn{সেঁকা}}, \emph{sẽkā}.) To steam? (\textbf{\bn{ভাপানো}}, \emph{bhāpāno}.) To sauté? (\textbf{\bn{সাঁতলানো}}, \emph{sā̃tlāno}.) Serving food? (\textbf{\bn{পরিবেশন}}, \emph{pôribeshôn}.) A mouthful? (\textbf{\bn{গ্রাস}}, \emph{grās}.) A flavour? (\textbf{\bn{স্বাদ}}, \emph{svād}.) Leftover from a meal? (\textbf{\bn{এঁটো}}, \emph{ẽṭo}.) To rinse after a meal? (\textbf{\bn{আঁচানো}}, \emph{ā̃chāno}.)
\end{tcolorbox}
